\section{自适应观测者：参数化读出、闭环更新与自指结构}

\subsection{观测者作为参数化读出系统}

设观测者由参数 $\theta$ 控制其读出协议：
\[
\theta := (E_{\parallel},W,\varepsilon,\ldots),
\]
其中包含投影子空间、接受窗、分辨率与可能的内部校准。观测数据记为
\[
y_t = \mathcal{R}_{\theta_t}\big(\Sigma(u_t)\big),
\]
其中 $\mathcal{R}_{\theta}$ 是从几何对象到离散记录的读出映射。

\subsection{相位平移与可控自由度}

在切投影模型中，余维相位偏移 $u$ 是决定 $\Sigma(u)$ 的关键参数。本文将“能动性”表述为：在不改变本体层 $\Lambda$ 与全局态的前提下，观测者可通过控制 $u$ 或 $\theta$ 改变其可观测切面与等价类分解，从而选择不同的可观测轨迹。形式上，这是参数空间上的控制变量，而非本体层结构的改变。

\subsection{闭环更新与自指}

为刻画“观测影响观测者自身协议选择”的自指结构，设更新律为
\[
\theta_{t+1} = \mathcal{U}(\theta_t, y_t),
\qquad
u_{t+1} = \mathcal{V}(u_t, y_t).
\]
该闭环系统在数学上是参数—数据耦合的离散动力系统。其“自指性”来自：$y_t$ 同时依赖于 $(\theta_t,u_t)$，而 $(\theta_{t+1},u_{t+1})$ 又由 $y_t$ 反馈决定。

\subsection{更新动因的协议化表述}

“为何更新”在本文中不以主观目的论表述，而以可审计的变分原则或约束满足表述。给定一个可计算的失配泛函
\[
\mathcal{J}(\theta,u; y_{0:t}) ,
\]
例如预测误差、约束违背预算、或信息增益的负值，更新律可取为
\[
(\theta_{t+1},u_{t+1})
=
\arg\min_{\theta,u}\ \mathcal{J}(\theta,u; y_{0:t}),
\quad \text{或其近似梯度步。}
\]
因此，观测者结构变化可被解释为：在有限资源与稳定性约束下，为降低失配而进行的参数重整。该表述不需要在本体层引入随机性或超出协议的目的性命题。

